
Large language models exhibit diverse failure modes---generating misinformation (truthfulness failures) or succumbing to adversarial perturbations (robustness failures)---yet these capabilities are typically studied independently. This work presents the first systematic correlation analysis between truthfulness and adversarial robustness across LLMs. Evaluating 14 decoder-only models from four families (Pythia, Llama-2, Mistral, Falcon) on TruthfulQA MC1 and AdvGLUE, a strong positive partial correlation (r=0.80, p<0.001, 95\% CI [0.08, 0.97]) emerges after controlling for model size. The natural hypothesis that model calibration underlies both capabilities is tested directly by measuring Expected Calibration Error (ECE) on MMLU. This mechanism is falsified: ECE shows no significant relationship with either metric (r=-0.12, p=0.68 for TruthfulQA; r=-0.16, p=0.58 for AdvGLUE), and poorly-calibrated models exhibit stronger truthfulness-robustness correlation than well-calibrated models---the opposite of what the calibration hypothesis predicts. These findings establish that truthfulness and robustness are linked through an unknown common cause.
